% ============================================================
% Chapter I — Introduction
% ============================================================
\section{Introduction}

\subsection{Background of the Study}

Formula One (F1) represents the pinnacle of motorsport engineering and strategy.
Each Grand Prix weekend produces terabytes of telemetry data---lap times, tire
degradation curves, fuel load estimates, weather conditions, and driver input
signals---that teams analyze in real time to gain competitive advantage. Despite
the sport's data-rich environment, predicting race outcomes remains an open
challenge because finishing position is determined by a complex interplay of
factors: raw car pace, tire compound selection, degradation management, fuel
burn-off, track architecture, and driver skill.

Historically, race outcome prediction has relied on either domain-expert
judgment or machine learning models trained on aggregate season statistics.
While these approaches achieve reasonable accuracy, they often treat each race
as an independent event and ignore the within-race physical dynamics---the
very variables that differentiate a podium finisher from a mid-field runner.

This study bridges that gap by engineering features directly from within-race
telemetry and applying multivariate statistical methods to model the
relationship between race strategy variables and a driver's finishing tier.
Rather than predicting exact finishing position (a continuous regression
problem with high variance), we frame the task as a three-class classification
problem: \textbf{Podium} (positions 1--3), \textbf{Point Scorer} (positions
4--10), and \textbf{No Points} (positions 11 and below). This tiered approach
reduces noise from unpredictable events---safety cars, first-lap incidents,
mechanical failures---while preserving the strategic information that
determines race outcomes.

\subsection{Conceptual Framework}

The study is grounded in the premise that a driver's finishing tier can be
modeled as a function of three categories of predictors:

\begin{enumerate}[label=(\arabic*)]
  \item \textbf{Baseline Pace Control} --- Grid Position captures the
    intrinsic car--driver combination's qualifying performance, absorbing
    the aerodynamic and engine capability differential between teams.
  \item \textbf{Tire Strategy} --- Tire Age (degradation) and Compound
    (Soft, Medium, Hard) represent the strategic choice of rubber and its
    wear trajectory over a stint.
  \item \textbf{Physical Interaction} --- Fuel Mass Proxy (linear burn-off
    from 110\,kg to 0\,kg) and the interaction term Degradation $\times$
    Fuel capture the combined effect of decreasing car mass and increasing
    tire wear---two forces that compound over a race distance.
\end{enumerate}

\noindent
The interaction term is mean-centered prior to multiplication to prevent
structural multicollinearity, allowing the coefficient to be interpreted as
the \textit{combined} effect of both forces rather than their individual
contributions.

\subsection{Statement of the Problem}

This study aims to determine whether within-race telemetry variables---tire
degradation, fuel burn-off, their interaction, and tire compound
selection---can statistically predict a driver's finishing tier when
controlling for baseline grid position. Specifically, it seeks to answer:

\begin{enumerate}
  \item Is there a statistically significant relationship between the
    engineered predictors and the probability of finishing in the Podium or
    Point Scorer tier relative to the No Points baseline?
  \item Does the interaction between tire degradation and fuel mass provide
    additional predictive power beyond their individual effects?
  \item Can a multinomial logistic regression model, trained on these
    predictors, generalize to unseen race circuits using leave-one-group-out
    cross-validation?
\end{enumerate}

\subsection{Hypotheses}

\begin{description}
  \item[$H_{0}$:] Tire degradation, fuel mass proxy, their interaction term,
    and tire compound selection have no statistically significant effect on
    the log-odds of a driver finishing in the Podium or Point Scorer tier,
    after controlling for grid position.
  \item[$H_{1}$:] At least one of the predictors---specifically the
    Degradation $\times$ Fuel interaction term---has a statistically
    significant effect on the log-odds of finishing in the Podium or Point
    Scorer tier ($p < 0.05$).
\end{description}

\subsection{Scope and Limitations}

This study uses lap-level telemetry data from the 2023 Bahrain, Saudi Arabian,
and Australian Grands Prix---covering permanent, street, and parkland circuit
types. The dataset comprises 2,880 observations after cleaning.

\textbf{Scope:}
\begin{itemize}
  \item The target variable is a three-class ordinal category (Podium, Point
    Scorer, No Points).
  \item Predictors are limited to grid position, tire age, fuel mass proxy,
    the degradation--fuel interaction term, and tire compound dummies.
  \item The statistical framework is multinomial logistic regression (MLE)
    validated via LOGO-CV.
\end{itemize}

\textbf{Limitations:}
\begin{itemize}
  \item The fuel mass proxy is an engineered linear approximation; actual fuel
    loads are proprietary and vary by team strategy.
  \item Three races may not capture the full variance of F1 outcomes across a
    23-race calendar with mixed weather, sprint formats, and variable
    DRS/regulation environments.
  \item DRS trains, safety car periods, and intra-team strategy calls are not
    explicitly modeled and contribute to residual variance in the mid-field.
\end{itemize}

\subsection{Significance of the Study}

This research contributes to the intersection of sports analytics and
multivariate statistical modeling by demonstrating that physically meaningful
telemetry features---not just aggregate season statistics---can predict race
outcomes with interpretable coefficients. The mean-centered interaction term
provides a template for modeling compound physical effects in other domains
where multiple degradation processes operate simultaneously. For the F1
analytics community, the LOGO-CV validation proves that the model captures
circuit-generalizable patterns rather than memorizing a single track layout.
